\documentclass[11pt]{article}
\usepackage[margin=2.5cm]{geometry}
\usepackage{amsmath}
\usepackage{booktabs}
\usepackage{parskip}

\title{Why a hazard model for incident plaque}
\date{}

\begin{document}
\maketitle

\section*{In one paragraph}

The model is still a binary regression: did a segment develop plaque between
the baseline and the follow-up scan, yes or no. ``Hazard model'' means only
that the link function, the complementary log-log (cloglog), comes from
describing plaque as something that can start at any moment during the
window at some rate, the hazard. It was not chosen because it fits better; a
logit link would fit the segment data about as well, and the likelihood-ratio
(LR) test of the geometry block works under either. It was chosen because of
how the outcome arises and what the model is used for afterwards.

\section*{Reason 1: segment risks combine into one person risk}

This is the reason that decided it. The question is about people, but the
model is fitted on nine segments per person. A person is free of plaque only
if every segment is. Under cloglog each segment has
$1-p_s = \exp(-e^{b}\Lambda_s)$, where $\Lambda_s$ is its cumulative hazard and
$b$ the person random intercept, so
\[
  P(\text{no plaque} \mid b) = \prod_{s=1}^{9} e^{-e^{b}\Lambda_s}
  = \exp\!\Big(-e^{b}\sum_{s}\Lambda_s\Big).
\]
The segment hazards simply add up to one person hazard, and the person risk
is one integral over $b$. A product of logistic terms has no such form. The
project record shows the switch: the earlier plan used a person-level
logistic regression, and cloglog became the main link when the pooled
segment model was adopted.

\section*{Reason 2: it matches how the outcome is produced and observed}

Plaque can start at any time during the roughly ten years, but it is only
seen at the follow-up scan. If risk accumulates at a rate $\lambda(t)$, the
chance of still being free of plaque at time $t$ is
$\exp\{-\Lambda(t)\}$, with $\Lambda(t)$ the hazard summed up to $t$. Under
proportional hazards, $\Lambda = \Lambda_0(t)\,e^{\eta}$, and
\[
  \log\{-\log(1-p)\} = \log\Lambda_0(t) + \eta .
\]
This is the cloglog link, derived rather than chosen (Prentice and Gloeckler
1978). Logit has no such derivation for this design.

\section*{Reason 3: follow-up length gets a meaning}

Follow-up varies (interquartile range 9.7 to 12.2 years). On the cloglog
scale, log follow-up enters with a coefficient $\gamma$, giving
$\Lambda_0(t) \propto t^{\gamma}$. A value of $\gamma = 1$ means a constant
hazard over time, and $\gamma > 1$ means the hazard rises over time. In the
current code $\gamma$ is estimated, not fixed at 1, so a rising baseline risk
is allowed for.

\section*{Reason 4: the same scale as SCORE2}

SCORE2 is a ten-year risk. Transformed by the cloglog link it becomes a log
cumulative hazard, so it enters the model on the same scale as the other
terms. The fitted SCORE2 slope is 0.722.

\section*{Reason 5: hazard ratios instead of odds ratios}

$\exp(\beta)$ is a hazard ratio per unit of a feature. Odds ratios overstate
effects when the outcome is common. For a hazard ratio of 1.5:

\begin{center}
\begin{tabular}{lcc}
\toprule
Risk in the reference group & Odds ratio & Level \\
\midrule
2\%  & 1.51 & segment \\
15\% & 1.56 & segment \\
42\% & 1.75 & person \\
\bottomrule
\end{tabular}
\end{center}

At segment level the two links nearly agree; they part at the person level,
where about 42\% developed plaque.

\section*{What if the risk is higher in the last years?}

Two cases need to be kept apart.

If everyone's risk rises over time, for example with ageing, only the
baseline hazard $\lambda_0(t)$ changes shape. This is not a breach of
proportional hazards, which only require the ratio between people to stay
constant. It is absorbed by the intercept and by $\gamma > 1$.

If instead the effect of geometry itself grows over time, hazards are not
proportional. The model still estimates a well-defined quantity, the ratio of
the cumulative hazards over the window,
\[
  e^{\beta} = \frac{\Lambda_1(t)}{\Lambda_0(t)}
  = \frac{\int_0^t \mathrm{HR}(s)\,\lambda_0(s)\,\mathrm{d}s}
         {\int_0^t \lambda_0(s)\,\mathrm{d}s},
\]
a weighted average of the hazard ratio over time, weighted by when plaque
develops. If most plaque develops late, the late years dominate. Only the
claim that the hazard ratio does not depend on the length of the window is
lost, and with follow-up varying by about $\pm$12\% the effect is small. The
two cases cannot be told apart, since one follow-up scan does not show when
the plaque started.

\section*{What does not depend on any of this}

The LR test of the geometry block, the change in C and the calibration of
the person risk compare fits of the binary ten-year outcome. They need no
assumption about when within the window the plaque started.

\section*{What to state in the thesis}

The hazard ratios rest on proportional hazards, which cannot be tested with
one follow-up scan; without them each is a weighted average over the window.
They are conditional on the person random intercept, not averaged over the
population, and they should not be compared with odds ratios from other
studies. One sentence carries the choice:

\begin{quote}
The cloglog link was used as the form a proportional-hazards model takes when
an event is known only to have occurred between two scans, and under it the
segment hazards of a person add to one person hazard, so the segment risks
combine into a single person risk.
\end{quote}

\end{document}
